\chapter{Participant Information Leaflet}
\label{app:pil}

\begingroup
\setlength{\parindent}{0pt}
\setlength{\parskip}{0.8em}

This appendix reproduces the Participant Information Leaflet used during recruitment and consent for the live study.

\textbf{Researcher:} Priyansh Nayak\\
\textbf{Student ID:} 25350660\\
\textbf{Programme:} MSc Computer Science - AVR\\
\textbf{School:} School of Computer Science \& Statistics, Trinity College Dublin\\
\textbf{Supervisor:} Gareth Young

\textbf{Invitation}

You are being invited to take part in a research study about the use of virtual reality for teaching computer architecture concepts. Before you decide whether to take part, it is important that you understand why the research is being carried out and what your participation would involve. Please read the following information carefully and ask any questions you may have before deciding.

\textbf{What is the purpose of the study?}

Computer architecture is often considered difficult to learn because many of its core processes are abstract and not directly visible. This study aims to evaluate whether a virtual reality learning environment can help support understanding of concepts such as instruction flow, pipelining, memory hierarchy, and stack behaviour. The study forms part of a Master's dissertation in Computer Science at Trinity College Dublin.

\textbf{Why have I been invited to take part?}

You are being invited because the study is seeking adult participants who are willing to take part in an evaluation of a virtual reality learning environment. Participants must be aged 18 or over and able to provide informed consent.

\textbf{Do I have to take part?}

No. Participation is entirely voluntary. If you decide to take part, you will be asked to sign a consent form before the study begins. You may withdraw at any time without penalty. If you choose to withdraw before your data has been anonymised and included in the analysis, your identifiable or linkable data will be deleted where possible.

\textbf{What will happen if I take part?}

If you agree to take part, you will attend one study session lasting approximately 45 to 60 minutes. During the session:

\begin{enumerate}
\item you will be briefed on the study and asked to provide written informed consent;
\item you will use a prototype virtual reality learning environment;
\item you will complete a small number of guided learning tasks;
\item you will complete short questionnaires on conceptual understanding, engagement, and usability;
\item you may also be invited to provide brief written open-ended feedback on your experience.
\end{enumerate}

\textbf{Are there any risks in taking part?}

There is a low risk of temporary physical discomfort from using the VR headset, including eye strain, dizziness, or motion sickness. If you feel any discomfort at any point, the session will stop immediately and you may remove the headset straight away. You may pause or withdraw from the study at any time without penalty.

\textbf{Are there any benefits to taking part?}

There is no direct benefit to you from taking part. However, you may gain some insight into the topic being studied, and your participation will contribute to research on virtual reality as a learning tool in computer science education.

\textbf{Will my taking part in this study be kept confidential?}

Any personal data collected for study administration, such as contact details used for arranging participation, will be kept separate from the research data. Research responses will be anonymised as early as practicable for analysis and reporting. Only the researcher and supervisor, where necessary, will have access to identifiable study administration data.

\textbf{What personal data will be collected?}

The study may collect limited personal data for administration purposes, such as your name and email address if needed for arranging participation. Research data will include questionnaire responses and written feedback. No special category personal data is intended to be collected for this study. The study uses a Meta VR headset provided to the researcher by the School of Computer Science \& Statistics. Participants will not be required to use their own Meta account. The research team will not intentionally collect Meta account data as part of the study. However, as third-party hardware and software are used, limited device- or platform-level processing outside the control of the research team may occur.

\textbf{How will my data be stored and protected?}

Personal data and study data will be stored on password-protected devices and university-approved storage systems. Identifiable administration data will be stored separately from research responses. Devices used for storage and access will be protected by password security, encryption, and up-to-date anti-virus protection.

\textbf{How long will my data be kept?}

Personal data held in identifiable or coded form, such as contact details used for study administration, will be retained for up to 12 months and then deleted. Research data will be anonymised as early as practicable for analysis and reporting.

\textbf{What will happen to the results of the study?}

The results of the study will be used for a Master's dissertation and may also be presented in academic work arising from the project. Any reporting of the findings will be done in anonymised form so that individual participants are not identified.

\textbf{Who is organising and funding the study?}

The study is being conducted by Priyansh Nayak as part of a Master's dissertation at Trinity College Dublin. The project is not externally funded.

\textbf{Who can I contact for further information?}

If you have any questions about the study, please contact:

Priyansh Nayak\\
25350660\\
MSc Computer Science - AVR

If you have concerns about this study and wish to contact someone independent of the research team, you may contact the relevant Trinity College Dublin research ethics process through the School of Computer Science \& Statistics.

\endgroup
